\documentclass[11pt]{article}
\usepackage[T1]{fontenc}
\usepackage{lmodern}
\usepackage[margin=1in]{geometry}
\usepackage{amsmath,amssymb,amsthm,mathtools}
\usepackage{booktabs,longtable,array}
\usepackage{microtype}
\usepackage{xcolor}
\usepackage{xurl}
\usepackage[colorlinks=true,linkcolor=blue!45!black,urlcolor=blue!45!black,citecolor=blue!45!black]{hyperref}
\hypersetup{pdftitle={Fixed-input Diophantine equations for the consistency of ZFC},pdfauthor={Research note prepared for ProveIt},pdfsubject={A 70-operation polynomial and an eight-equation system},pdfkeywords={ZFC, consistency, Diophantine equations, K215, Pell equations}}
\usepackage{enumitem}
\usepackage{fancyhdr}
\usepackage{listings}
\setlist{itemsep=3pt,topsep=5pt}
\setlength{\headheight}{14pt}
\pagestyle{fancy}
\fancyhf{}
\lhead{Diophantine consistency of ZFC}
\rhead{ProveIt research note}
\cfoot{\thepage}
\newtheorem{theorem}{Theorem}[section]
\newtheorem{lemma}[theorem]{Lemma}
\newtheorem{proposition}[theorem]{Proposition}
\newtheorem{corollary}[theorem]{Corollary}
\theoremstyle{definition}
\newtheorem{definition}[theorem]{Definition}
\newtheorem{remark}[theorem]{Remark}
\newcommand{\N}{\mathbb N}
\newcommand{\Z}{\mathbb Z}
\newcommand{\Pos}{\mathbb Z_{>0}}
\newcommand{\Con}{\operatorname{Con}}
\newcommand{\Halt}{\operatorname{Halt}}
\newcommand{\pop}{\operatorname{pop}}
\newcommand{\bits}{\operatorname{bits}}
\newcommand{\repo}{\texttt{cdd582bec13f908df3711fa9e05d1ddac368a5db}}
\newcommand{\norms}{\mathcal P}
\lstset{basicstyle=\ttfamily\small,breaklines=true,columns=fullflexible,keepspaces=true,frame=single,rulecolor=\color{black!20}}
\title{Fixed-input Diophantine equations for the consistency of ZFC\\[5pt]
\large A 70-operation polynomial and an eight-equation system}
\author{Research note prepared for the ProveIt repository}
\date{22 October 2026}

\begin{document}
\maketitle
\begin{abstract}
We specify an ordinary integer-coefficient polynomial in sixteen strictly
positive unknowns whose absence of solutions is provably equivalent in ZFC
to the syntactic consistency of ZFC. Its evaluation uses 70 binary arithmetic
operations, comprising 40 multiplications and 30 additions or subtractions,
when its five fixed integer coefficients are supplied as constants. Its exact
total degree is 169. A 71-operation variant has degree 156. A more transparent
alternative is an eight-equation system in seventeen positive unknowns,
evaluated by 61 arithmetic operations; its sum of squared residuals has degree
56 and an 84-operation evaluator.

The main simplification is specific to a fixed input. The repository's
universal ordinary-input construction needs a separate Pell block to encode
the input marker. Here that marker is a fixed coefficient, so the block can
be removed. A second reduction permits a signed repunit condition and proves
its sign only after recovering a power of two. Reusing this not-yet-verified
repunit saves a further operation when the first scale is changed from
$wq$ to $wq^2$.

The coefficients are fully determined by the literal K215 transition table,
an explicit 435-state normalization, a finite window predicate, and a finite
arithmetic recipe. Their enormous binary expansions and the complete window
list are not materialized. This is an ordinary mathematical proof, with
reproducible circuit and finite-controller checks, building on the pinned
repository's Pell, compiler, and consistency-calibration theorems. It is not
a newly kernel-checked end-to-end theorem, an internal ZFC proof object,
or a claim of global minimality.
\end{abstract}

\clearpage
\tableofcontents
\newpage

\section{The precise result and the meaning of simplicity}

Throughout, $\N=\{0,1,2,\ldots\}$ and $\Pos=\{1,2,\ldots\}$.
Let $\Con(\mathrm{ZFC})$ mean that the fixed ordinary first-order axiom system
for ZFC has no finite derivation of falsity. The derivability convention is
the one calibrated by the repository's K215 endpoint. Equivalent effective
presentations of the same usual theory can be substituted using the usual
formalized proof translations.

\begin{theorem}[Main result]\label{thm:main}
The coefficient recipe in Section~\ref{sec:coefficients} determines five
specific positive integers $\beta,K,M,N,W$. With these constants, the
polynomial $\norms_{70}$ defined in Section~\ref{sec:polynomial} satisfies
\[
 \mathrm{ZFC}\ \vdash\quad
 \bigl(\exists\mathbf x\in\Pos^{16}\quad
       \norms_{70}(\mathbf x)=0\bigr)
 \ \longleftrightarrow\ \neg\Con(\mathrm{ZFC}).
\]
Consequently ZFC proves
\[
 \boxed{\quad
 \bigl(\nexists\mathbf x\in\Pos^{16}\quad
       \norms_{70}(\mathbf x)=0\bigr)
 \ \longleftrightarrow\ \Con(\mathrm{ZFC}).\quad}
\]
No consistency premise is used to prove either equivalence.
\end{theorem}

The reference to ZFC provability means that the mathematical argument can be
carried out in ordinary ZFC. It does not assert that this archive contains a
finite string in a chosen ZFC proof language that has been mechanically
checked as that derivation. Section~\ref{sec:logical-status} separates this
claim from the repository's external Lean and Rocq theorems.

There is no invariant, universally agreed meaning of ``the simplest
Diophantine equation.'' Degree, variable count, coefficient height, expanded
monomial count, arithmetic-circuit size, and proof length are different
measures. This note makes its principal comparison using the repository's
arithmetic model: each binary $+$, $-$, or $\times$ costs one operation;
reuse of an already computed expression and access to a fixed integer cost
zero. Multiplication by a fixed coefficient still costs one multiplication.
Equality tests are listed separately for a system. The finalizer is charged
when a system becomes one polynomial.

\begin{center}
\begin{tabular}{lrrrr}
\toprule
Construction & Positive unknowns & Equations & Operations & Degree\\
\midrule
Main polynomial $\norms_{70}$ & 16 & 1 & $40M+30A=70$ & 169\\
Lower-degree variant $\norms_{71}$ & 16 & 1 & $40M+31A=71$ & 156\\
Budget-square variant $\norms_{73}$ & 16 & 1 & $41M+32A=73$ & 157\\
Transparent system & 17 & 8 & $33M+28A=61$ & $\leq28$\\
Its sum of squares & 17 & 1 & $41M+43A=84$ & 56\\
Inherited universal polynomial & $18+\text{input}$ & 1 & $47M+37A=84$ & 187\\
\bottomrule
\end{tabular}
\end{center}

The last row concerns a uniform positive input and has a different contract.
Our construction specializes the input and adds its marker as a fixed
coefficient. Thus 70 is a fixed-instance improvement, not a new uniform
universal-operation bound. None of these numbers is a lower bound over all
possible Diophantine descriptions. The very large coefficient height is
reported explicitly rather than included implicitly in the operation count.

\section{The polynomial, with every variable and coefficient specified}
\label{sec:polynomial}

The sixteen unknowns are
\[
 \boxed{J,F,\alpha,t,f,h,i,T,s,w,\tau,\eta,\zeta,y,Z,\gamma\in\Pos.}
\]
The letter $T$ is an auxiliary quotient and $t$ is a transport quotient.
The uppercase $Z$ is a packed word; it is unrelated to the ring $\Z$.
All expressions introduced next are polynomial abbreviations, not additional
unknowns or uncharged arithmetic primitives.

First define the geometry and scales:
\begin{align}
 C&=Z+W,& u&=C+\alpha,& q&=u+F,\label{eq:geometry}\\
 v&=u-Z,& P&=\beta J,& X&=wq^2,\nonumber\\
 Y&=sq^3,& E&=XY,& k&=\eta+\zeta,\label{eq:scales}\\
 c&=kY+\eta,& a&=E+Y,& H&=4a+3,\nonumber\\
 \Delta&=a^2+H=(a+1)(a+3),&
 D&=X+ac+\gamma H,& L&=E(kY).\label{eq:pell-data}
\end{align}
Next define the packed index and auxiliary expressions:
\begin{align}
 R&=(Pu+v)(q^2-1)+(M+qN)J,\label{eq:packed70}\\
 V&=c(Tf-1)-Rf^2,\label{eq:bezout}\\
 A_*&=(i\Delta c^2)^2.\label{eq:aux-coeff}
\end{align}
The seven factors are
\begin{align}
 U_0&=q-P,\label{eq:factors}\\
 U_1&=\tau^2-L(L+k),\nonumber\\
 U_2&=D^2-\Delta c^2,\nonumber\\
 U_3&=A_*(V^2-y^2)+y^2,\nonumber\\
 U_4&=k-hE-R,\nonumber\\
 U_5&=(K+w)C+u-tP,\nonumber\\
 U_6&=\Delta f^2-A_* .\nonumber
\end{align}
The promised equation is
\begin{equation}\label{eq:main-polynomial}
 \boxed{\qquad
 \norms_{70}=U_0U_1U_2U_3U_4U_5U_6-\Delta=0.
 \qquad}
\end{equation}

This is one ordinary polynomial equation with integer coefficients. All
exponents of unknowns in its definition are fixed small integers. In
particular, the coefficient $W=2^{b+2d}$ introduced below is a fixed integer;
there is no existential exponent, computation-length parameter, hidden
bounded universal quantifier, or oracle in~\eqref{eq:main-polynomial}.
Appendix~\ref{app:circuit} gives the complete 70-row straight-line evaluator.

\subsection{The five coefficients}

The exact compiler recipe produces integers $b,L_0,d,B,D_C,D_R,M_C,M_F$,
where $d=bL_0$ and $B=2^d$. The five ports used by the polynomial are
\begin{equation}\label{eq:five-coefficients}
 \boxed{
 \beta=B-1,\quad K=D_C+BD_R,\quad M=M_C,\quad
 N=M_F+B-1,\quad W=2^{b+2d}.}
\end{equation}
The $M_F$ in this formula is the native verification mask. The polynomial's
$N$ is its \emph{shifted} source value. Omitting this shift gives a different
and unjustified equation.

For the explicit K215 instance supplied here,
\[
 b=5^{44},\qquad L_0=5^{45},\qquad d=5^{89},
 \qquad B=2^{5^{89}},
 \qquad W=2^{5^{44}+2\cdot5^{89}}.
\]
The other coefficients are fixed finite sums of powers of $2^b$, whose
terms are specified by the literal controller and window predicate. These
are exact descriptions of integers, not parameters waiting to be selected
by a halting or consistency decision. Section~\ref{sec:coefficients} and
\texttt{artifacts/compiler/COEFFICIENTS.md} give the full recipe.

\section{What is inherited, and what changes here}

The general connection between computably enumerable sets and Diophantine
equations is classical. Jones's explicit universal constructions and the
DPRM theorem provide the historical framework~\cite{Jones1980,Jones1982}.
Explicit provability encodings also predate this work~\cite{CarlMoroz},
and full formal DPRM developments are available~\cite{FormalDPRM}.
The contribution here is a specific reduction and instantiation of the
much more recent constructions already present in ProveIt.

All repository dependencies are pinned at \repo. The argument uses three
substantive inherited components.

\begin{enumerate}
\item \textbf{Consistency calibration.} The exact K215 machine halts from
blank tape exactly when there is a finite ZFC refutation. The ordinary
proof says this equivalence is formalizable in ZFC; external Lean and Rocq
proof lanes establish their corresponding endpoint statements~\cite{K215}.
\item \textbf{Pell decoding.} The strong half-binomial kernel, its auxiliary
rank and signed step-down lemmas, and the asymmetric-scale argument recover
the actual main Pell index and a power of two from polynomial norm and
ratio conditions~\cite{Kernel,Asymmetric,Bezout}.
\item \textbf{The fixed helical compiler.} The modified75 layout turns
finite local windows into fixed coefficients, with shifted masks,
synchronization, a unique start/end marker argument, and a five-adic dummy
adjustment for completeness~\cite{Compiler75,CompilerDetails}.
\end{enumerate}

We do not replace any of these by numerical samples. The new argument
removes the input-Pell block, proves a weaker initial repunit condition
sufficient, reuses its multiple before the sign is known, and supplies an
explicit controller instance. The proof below records precisely where the
inherited local lemmas are used. It does not invoke the whole universal
polynomial theorem with one of its hypotheses silently deleted.

\section{Why the input-Pell block can be deleted}
\label{sec:fixed-input}

In the modified75 compiler, ordinary input $x>0$ is represented by a single
End bit at the integer
\[
 W_x=2^{2dx+b}.
\]
A four-equation Pell block is needed to enforce this expression while $x$
varies. We fix ordinary input $x_*=1$, so the marker is the coefficient
$W=W_{x_*}$ in~\eqref{eq:five-coefficients}. The identity $C=Z+W$ inserts
that marker directly. No variable exponent remains to be enforced.

\begin{lemma}[Fixed-marker specialization]\label{lem:fixed-marker}
For coefficients from the modified75 compiler, retain its full native
kernel, packing, and transport conditions, set $C=Z+2^{b+2d}$, and retain a
positive bound $C<q$. Its soundness argument decodes ordinary input $1$
without the four input-Pell equations. Every accepting computation at input
$1$ still has positive witnesses for these retained conditions.
\end{lemma}

\begin{proof}
The native kernel and packing argument recover $q=B^m$ and the shifted mask
conditions. Since $W<C<q$ and $2^b<B$, necessarily $m>2$.
The native $Z$ word has its End selector forbidden, so adding $W=2^bB^2$
has no carry and inserts exactly one End at cell distance $2$. The shifted
low mask supplies Start at the origin. The existing marker bijection then
makes Start unique and identifies these two markers with the same
initialization run. The run has three I cells and its Q head, hence four
unary $1$ symbols. This is ordinary input $1$ in the modified75 convention.
Conversely, every old complete witness has exactly this marker; forgetting
its five private input-Pell coordinates and treating $W$ as a coefficient
preserves all retained equations. The canonical padding theorem supplies
any stronger positive budget used below.
\end{proof}

The explicit normalizer in Section~\ref{sec:controller} erases its unary
initialization and runs K215 on a blank tape. It therefore has the same
halting behavior for every allowed unary length. Nevertheless, the marker
proof is still needed: it guarantees a genuine initialized computation
rather than an unrelated periodic component.

\section{Soundness of the 70-operation polynomial}
\label{sec:soundness}

Fix the five actual coefficients and a positive tuple satisfying
\eqref{eq:main-polynomial}. We show that the tuple decodes an accepting
computation. This direction is where false solutions must be excluded.

\subsection{Cancel the positive discriminant before treating factors as units}

All definitions of $a$ are positive, so
$\Delta=(a+1)(a+3)>0$ before imposing any equation. Define
\[
 S=f^2-\Delta i^2c^4.
\]
Then $U_6=\Delta S$, and the exact polynomial identity is
\begin{equation}\label{eq:cancel}
 \norms_{70}
 =\Delta\bigl(U_0U_1U_2U_3U_4U_5S-1\bigr).
\end{equation}
Consequently the seven integer factors inside the parentheses have product
$1$, and each is $1$ or $-1$. Inferring unit factors directly from a product
equal to $\Delta$ would be invalid; the cancellation in~\eqref{eq:cancel}
is essential.

\begin{lemma}[The four norm signs]\label{lem:norm-signs}
At a positive zero, $U_1=U_2=U_3=S=1$.
\end{lemma}

\begin{proof}
Put $V_0=XY^2$. The identity
$L(L+k)=V_0(V_0+1)k^2$ shows that $U_1$ is a Pell norm.
For $V_0>1$, the equation
$r^2-V_0(V_0+1)\ell^2=-1$ has no solution with $\ell>0$.
Indeed, such a solution gives
$V_0\ell<|r|<(V_0+\tfrac12)\ell$ and hence the smaller positive integer
\[
 \ell'=(2V_0+1)\ell-2|r|,\qquad 0<\ell'<\ell.
\]
Multiplying by the norm-one unit
$2V_0+1-2\sqrt{V_0(V_0+1)}$ produces another norm-minus-one solution
with this coefficient, contradicting minimality. Thus $U_1=1$.

The discriminant $\Delta=a^2+4a+3$ is $0$ or $3$ modulo $4$.
A norm $r^2-\Delta z^2$ therefore cannot be $-1$ modulo $4$.
This proves $U_2=S=1$. Finally, $A_*$ is a square. Modulo $4$, the expression
$A_*V^2-(A_*-1)y^2$ is a square whether $A_*$ is $0$ or $1$ modulo $4$.
It cannot be $-1$, so $U_3=1$ as well.
\end{proof}

Write
\[
 U_0=\epsilon,\qquad U_4=\iota,\qquad U_5=\nu,
 \qquad\epsilon,\iota,\nu\in\{1,-1\},\quad
 \epsilon\iota\nu=1.
\]
We will establish their positive signs in the order $\iota$, $\epsilon$,
$\nu$. In particular the exact repunit relation is not available yet.

\subsection{Bounds valid for both repunit signs}

From the definitions, $C=Z+W>0$, $q=C+F+\alpha$, and $F+Z<q$.
Also $q>W\geq2B^2$, so certainly $q\geq16$. Since
$\epsilon=q-\beta J$, we have $\beta J=q-\epsilon$ and
\begin{equation}\label{eq:j-bound}
 J\leq(q+1)/\beta.
\end{equation}
The compiler masks satisfy
\[
 0<M<\beta,\qquad 0<M_F<\beta,
 \qquad N=M_F+\beta<2\beta.
\]
Let $G=q^2-qF-Z$. The positive budget gives
\[
 2q-1\leq G\leq q^2-q-1.
\]
Crucially, the actually computed first factor in~\eqref{eq:packed70} is
\begin{equation}\label{eq:gap-error}
 Pu+v
 =G+(1-\epsilon)(q-F).
\end{equation}
It is not permissible to replace it by $G$ before proving $\epsilon=1$.
Since $1\leq q-F\leq q-1$, both signs imply
\[
 2q-1\leq Pu+v\leq q^2+q-3.
\]
Using~\eqref{eq:j-bound} in the mask term yields
\begin{align}
 R&>(2q-1)(q^2-1)>3q+1,\label{eq:r-lower}\\
 R&<(q^2+q-3)(q^2-1)+(1+2q)(q+1)\nonumber\\
  &=q^4+q^3-2q^2+2q+4<q^4+q^3.\label{eq:r-upper}
\end{align}
The revised scale $X=wq^2$, $Y=sq^3$ now supplies
\begin{equation}\label{eq:margin}
 E\geq q^5>R+5,
 \qquad a=E+Y>R+2.
\end{equation}
This extra margin is exactly why the repunit can be reused while its sign
is still unknown. The former asymmetric scale $X=wq$ would not supply this
bound for the negative branch of~\eqref{eq:gap-error}.

\subsection{Restore the auxiliary congruence witnesses}

The equality $S=1$ gives
\begin{equation}\label{eq:strong}
 f^2=1+\Delta i^2c^4.
\end{equation}
We have $c>2$, hence $f>c^2>2c$. Since $R<a<\Delta$,
\begin{align}
 A_*-(Rf^2+c)
 &=(\Delta-R)f^2-\Delta-c\nonumber\\
 &\geq f^2-\Delta-c
 =1+\Delta(i^2c^4-1)-c>0.\label{eq:aux-margin}
\end{align}
The auxiliary norm is
\[
 A_*V^2-(A_*-1)y^2=1.
\]
It excludes $V=0$. If $v_0=|V|>1$, then
$v_0^2-1=(A_*-1)(y^2-v_0^2)$ and $y\geq v_0+1$.
It follows that $v_0\geq2A_*-1$. But
$V=c(Tf-1)-Rf^2>-Rf^2-c$, so~\eqref{eq:aux-margin}
excludes every negative value with $|V|>1$.
The remaining possibilities $V=\pm1$ are impossible because
$V\equiv-c\pmod f$ and $2<c<f/2$. Thus $V>0$.

Define the proof-only integers
\begin{equation}\label{eq:restore-jo}
 o=cT-Rf,
 \qquad j=Tf-R\Delta i^2c^3-1.
\end{equation}
Then $of=V+c>0$, and~\eqref{eq:strong} gives $jc=V+R>0$.
Thus $o,j>0$ and
\begin{equation}\label{eq:aux-congruences}
 V=of-c=jc-R.
\end{equation}
The positivity is proved before applying a Pell-index theorem; it is not
an assumption about the omitted coordinates.

\subsection{The actual Pell index and the sign of the index factor}

Write $A=a+2$ and $P_0=2XY^2+1$. For $A\geq2$ use
\[
 \chi_A(n)+\psi_A(n)\sqrt{A^2-1}
 =(A+\sqrt{A^2-1})^n.
\]
The first and main positive norms give positive indices $n,p$ with
\[
 k=2\psi_{P_0}(n),\qquad
 \tau=\chi_{P_0}(n),\qquad
 c=\psi_A(p),\qquad D=\chi_A(p).
\]
Here $P_0>A$, $c>k$, and the two ratio slacks give
\begin{equation}\label{eq:ratio}
 kY<c<k(Y+1).
\end{equation}
Parameter monotonicity implies $p>n$. Since $P_0\equiv1\pmod E$,
\[
 2n\equiv R+\iota\pmod E,
 \qquad n\geq(R-1)/2.
\]
The lower bound on $R$ makes $p\geq25$. Standard Pell growth now gives
$c>A\Delta^2$ and $c>2p$. Also
$k=R+\iota+hE$ and $Y\geq q^3$ give $c>2R$.

At this point the inherited strong-rank and signed auxiliary step-down
lemmas apply~\cite{Kernel,Bezout}. To match their ordinary strong premise,
put $i_{\mathrm{old}}=\Delta i$; then
\[
 (i_{\mathrm{old}}c^2)^2=\Delta(f^2-1).
\]
Their other hypotheses are the just-established positive $o,j,V$, the
congruences~\eqref{eq:aux-congruences}, $c>A\Delta^2$, $c>2p$, and $c>2R$.
They identify the main index as
\begin{equation}\label{eq:p-R}
 p=R.
\end{equation}
These are local applications of the rank and step-down lemmas. They do not
use a decoded input, a typed word, the exact repunit equation, or the final
universal-polynomial theorem.

The least positive representative of $R+\iota$ modulo $E$ is itself.
If $2n=R+\iota+vE$ with $v\geq1$,~\eqref{eq:margin} would give
$2n>2R$, contradicting $n<p=R$. Hence $2n=R+\iota$.
If $\iota=-1$, then $R=2n+1$. Put $Q_2=\chi_A(2)=2A^2-1$.
Since $Q_2>P_0$ and $A>Y+1$, Pell duplication gives
\[
 \psi_A(2n)=2A\psi_{Q_2}(n)\geq A k,
 \qquad
 c=\psi_A(2n+1)>A k>k(Y+1),
\]
contradicting~\eqref{eq:ratio}. Therefore
\begin{equation}\label{eq:index-positive}
 \iota=1,\qquad 2n=R+1.
\end{equation}

\subsection{Recover a power of two before deciding the repunit sign}

Set $r=n-1$, so $R=2r+1$ and $r\geq1$. The lower Pell-ratio estimate used
in the asymmetric-scale proof~\cite{Asymmetric} is
\[
 \frac{c}{k/2}>
 \xi,\qquad \xi=\frac{(X+1)^{2r}}{X^r}.
\]
Its elementary hypotheses $X\geq16$, $Y\geq4096$, and $6XY^2>a$ hold
here. This lower estimate does not require a decoded word or an upper
binomial-floor error estimate. Combining it with~\eqref{eq:ratio} yields
\[
 Y>\xi/2-1>X^r/3,
 \qquad a=Y(X+1)>X^{r+1}/3>2\cdot4^r=2^R.
\]
The sequence $\chi_A(j)-a\psi_A(j)$ is congruent to $2^j$ modulo
$H=4a+3$, as follows directly from its recurrence and initial values $1,2$.
The definition of $D$ therefore gives $X\equiv2^R\pmod H$.
Both numbers lie strictly between $0$ and $a$, so
\begin{equation}\label{eq:dyadic}
 X=2^R.
\end{equation}
Since $q\mid X$, we have $q=2^\ell$ for some $\ell\geq4$.

\begin{lemma}[The negative repunit sign is impossible]\label{lem:rep-sign}
Let $B=2^d$ with $d\geq3$. If $q$ is a positive power of two and
$q-(B-1)J\in\{1,-1\}$ for a positive integer $J$, then
$q-(B-1)J=1$.
\end{lemma}

\begin{proof}
For the negative sign, $2^\ell\equiv-1\pmod{2^d-1}$.
Writing $\ell=ad+r_0$ with $0\leq r_0<d$ gives divisibility of
$2^{r_0}+1$ by $2^d-1$. This is impossible because
\[
 0<2^{r_0}+1\leq2^{d-1}+1<2^d-1.
\]
\end{proof}

Thus $\epsilon=1$, so $P=q-1$. The product of the remaining unit signs
now forces $\nu=1$. Moreover $q\equiv1\pmod{2^d-1}$ implies $d\mid\ell$
by the same reduction of the exponent; hence $q=B^m$.

\subsection{Restore the native compiler equations and decode the computation}

All preliminary errors have disappeared:~\eqref{eq:gap-error} becomes
$Pu+v=q^2-qF-Z$, and the packing is exactly
\[
 R=(q^2-Z-qF)(q^2-1)+(M_C+q(M_F+B-1))J.
\]
The original mask bounds now give $3q+1<R<q^4$.
Equation~\eqref{eq:dyadic} and $R>3q>3\ell$ imply $q^3\mid X$, so
\[
 w_{\mathrm{old}}=X/q^3=w/q\in\Pos.
\]
The proof-only substitution $i_{\mathrm{old}}=\Delta i$ and the positive
$o,j$ from~\eqref{eq:restore-jo} restore the ordinary strong half-binomial
kernel. Every restored definition and norm agrees exactly. If the preceding
projection presentation keeps two positive slacks, they can also be restored:
$c>A\Delta^2$ and the main norm imply $D>(a+1)c$, so
$\gamma H=D-ac-X>c-X>A\Delta^2-a>H$.
Thus $\gamma\geq2$; take $\rho=1$ and
$\sigma=\gamma-1$. No extra equation is needed.

The transport condition $U_5=1$ reads
\[
 (K+w)C-F=(t-1)(q-1).
\]
Consequently the positive integer
\[
 z=t-1+w(q+1)C
\]
satisfies
\[
 (K+X)C=F+z(q-1).
\]
This is the native transport equation since $X=wq^2$.
Also $C<q$, $Z<C$, and $C=Z+W$ hold by definition.

The inherited half-binomial population argument and modified75
synchronization now apply without alteration. They yield the typed cyclic
word, the actual whole-cell temporal shift, and the local transition tests.
Lemma~\ref{lem:fixed-marker} supplies the unique initialization markers at
ordinary input $1$. The finite helical soundness theorem yields a halting
run of the explicit normalizer, hence a blank halting run of K215. This
proves soundness.

\section{Completeness and the ZFC equivalence}
\label{sec:completeness}

Suppose K215 halts on blank tape. The explicit normalized controller then
halts on four unary $1$ symbols. The inherited helical padding and modified75
compiler completeness theorems produce a positive native witness at ordinary
input $1$. The normalized strong and auxiliary-quotient refinements used in
the repository's 84-operation polynomial supply all of the corresponding
positive coordinates~\cite{Universal84,Bezout}.

Keep its $q,C,F,Z,J,k,c,a,D,f,h,i,T,s,\tau,\eta,\zeta,y$ and its main
power $X$. Define the new positive quantities
\[
 W=2^{b+2d},\qquad
 \gamma=\rho+\sigma,\qquad
 \alpha=q-C-F,\qquad
 w=X/q^2.
\]
The old stronger budget is
$q=C+F+Z+\alpha_{\mathrm{old}}+2d$, so
$\alpha=Z+\alpha_{\mathrm{old}}+2d>0$.
The old main power is divisible by $q^3$, hence the new $w$ is a positive
integer. The old input theorem gives exactly $C=Z+W$.

The native transport equation implies that
\[
 t=1+\frac{(K+w)C-F}{q-1}
\]
is an integer: $X\equiv w\pmod{q-1}$. This quotient is greater than $-1$
because $0<F<q$ and $(K+w)C>0$. Therefore $t\geq1$.
This positivity check is necessary when changing the scale and the transport
quotient; it does not follow merely from the positivity of the old quotient.

The new definitions recover every retained native expression. The factors
$U_0,U_1,U_2,U_3,U_4,U_5$ equal $1$ and $U_6=\Delta$, so
$\norms_{70}=0$. This proves completeness.

Combining the two directions gives
\[
 \exists\mathbf x\in\Pos^{16}\ \norms_{70}(\mathbf x)=0
 \quad\longleftrightarrow\quad
 \Halt(K215,\mathrm{blank}).
\]
The inherited calibration
$\Halt(K215,\mathrm{blank})\leftrightarrow\neg\Con(\mathrm{ZFC})$
then proves Theorem~\ref{thm:main}. Every step is ordinary finite arithmetic,
Pell theory, finite-computation coding, or an explicitly cited set-theoretic
calibration. These arguments can be formalized in ZFC.

The construction does not prove that either side is false. If ZFC is
consistent, the equation has no positive solution, and ZFC cannot prove
that absence under the usual hypotheses of the second incompleteness
theorem. Mere consistency is not being used to assert that ZFC also cannot
prove the existence of a solution; that stronger independence statement
requires an appropriate soundness hypothesis. Only the displayed
equivalence is needed here.

\section{An eight-equation alternative}
\label{sec:system}

For a more direct audit, use the same five coefficients and seventeen
positive unknowns
\[
 J,F,\alpha,z,Z,f,h,i,j,o,s,w,\tau,\eta,\zeta,\gamma,y.
\]
Define
\begin{align*}
 C&=Z+W,& q&=C+F+\alpha,& X&=wq^3,&Y&=sq^3,\\
 E&=XY,& k&=\eta+\zeta,&c&=kY+\eta,&a&=E+Y,\\
 H&=4a+3,&\Delta&=a^2+H,&D&=X+ac+\gamma H,&L&=E(kY),\\
 R&=(q^2-Z-qF)(q^2-1)+(M+qN)J,\\
 A_0&=\Delta(f^2-1),&V&=of-c,&U&=jc-R.
\end{align*}
The eight equations are
\begin{equation}\label{eq:system8}
\boxed{\begin{aligned}
 \beta J&=q-1,\\
 (K+X)C&=F+z(q-1),\\
 \tau^2-L(L+k)&=1,\\
 k-hE-R&=1,\\
 D^2-\Delta c^2&=1,\\
 (ic^2)^2&=A_0,\\
 A_0(V^2-y^2)+y^2&=1,\\
 U&=V.
\end{aligned}}
\end{equation}

Here $R$ is a computed expression, not a positive supplied coordinate.
The budget $q=C+F+\alpha$ guarantees $F+Z<q$ and therefore the required
positive range of $R$. Dropping $F$ from this budget while also eliminating
the supplied $R$ would leave a gap in the proof.

The first norm identity
\[
 (E^2+X)(kY)^2=L(L+k)
\]
holds as a polynomial identity. The strong equation makes
$A_0=(ic^2)^2$ at a simultaneous zero, and $U=V$ restores the original
auxiliary argument. Thus~\eqref{eq:system8} is exactly a fixed-marker
projection of the ordinary symmetric half-binomial construction. The
previous soundness and completeness proof applies, without the signed
repunit optimization or the positive Bezout quotient elimination.

The emitted comparison circuit costs 61 operations. If $r_1,\ldots,r_8$
are its literal left-minus-right residuals, then
\begin{equation}\label{eq:sos}
 \boxed{r_1^2+\cdots+r_8^2=0}
\end{equation}
is an equivalent single polynomial equation over either the positive or
nonnegative reparameterized domain. Its finalizer costs eight subtractions,
eight squares, and seven additions, giving $61+23=84$ operations.

The residual degree bounds, in the displayed order, are
\[
 1,5,26,9,22,22,28,6.
\]
The seventh has the unique degree-28 leading part
$w^2s^4k^2q^{18}f^2$, where the leading part of $q$ is
$Z+F+\alpha$. Hence the sum of squares has exact degree 56.
The auxiliary coefficient $A_0=\Delta(f^2-1)$ and argument $V=of-c$
are intentional: replacing them off the zero set by their equivalent
high-degree expressions would worsen the polynomial degree without changing
the simultaneous solutions.

\section{The two nearby single-polynomial variants}

The archive also contains two useful checkpoints.

For $\norms_{71}$ keep the same positive variables and geometry
$C=Z+W$, $q=C+F+\alpha$, but set $X=wq$ and use the exact expression
$q-1$ in packing and transport. Keep the signed factor
$q-\beta J$. Its initial packed-index bound is
\[
 R<q^4-q^3+4q+2,
\]
which already gives $E\geq q^4>R+5$. The same local argument recovers the
dyadic $q$ and then the positive repunit sign. Its evaluator has 71 rows
and exact degree 156. It trades one additional subtraction for a smaller
degree than the main construction.

For $\norms_{73}$ impose $q=\beta J+1$ by definition, set $C=Z+W$,
take $X=wq$, and replace the signed repunit factor by
\[
 1+(q-F-C-\alpha)^2.
\]
At a zero this positive unit is $1$, so the budget holds exactly.
All other factors are the same normalized native expressions. This variant
avoids the delayed repunit-sign proof; it has 73 operations and exact
degree 157. The 70, 71, and 73 circuits represent the same acceptance
condition but are not asserted to be identical polynomials off their zero
sets or to have identical positive witness tuples.

\section{An explicit finite program and coefficient recipe}
\label{sec:coefficients}

\subsection{Pinning the actual machine}
\label{sec:controller}

The input is the exact released dense table
\texttt{candidate-k215/k215.tm}, with SHA-256
\begin{center}\small\ttfamily
 fdbcef0825886035c3091c8537a69535bc96f9e265f0ae45863543d40ae70090.
\end{center}
Its working states are $0,\ldots,214$, its separate halt is $215$, and its
tape symbols are $0,1$. It starts in state $0$ at position $0$ on the blank
two-way tape. The archive includes this table and its upstream attribution.

The normalization is completely explicit. It has three loader states,
one nonhalting trap, one absorbing halt, and two copies of each of the 215
working states, for 435 states in total. At physical position $-j$,
$j\geq0$, a pair stores original cells $j$ and $-j-1$. A side bit in the
state selects the component being read. Four ordinary pair symbols and four
origin-flagged pair symbols, together with loader symbols $1,2$, give ten
tape symbols. Blank is the ordinary pair $(0,0)$.

Away from the flagged origin, moving further out on either side of the
original tape moves left physically, while moving toward the origin moves
right. At the flagged origin, crossing between original cells $0$ and $-1$
changes the side tag and stays in place. The written component changes
exactly as in K215 and the other component remains unchanged. The six-case
motion table and every one of the 4,350 transitions are supplied in
\texttt{artifacts/compiler/}.

The first transition, from the unary initialization, writes $2$ and stays.
The loader scans left erasing the remaining unary $1$s, returns right across
the erased zeros, replaces the $2$ by a flagged blank pair, and enters the
copy of original state $0$ on side zero. It never visits a positive physical
position. On the required four unary $1$s this takes ten steps. Induction on
the original transitions proves exact tape and state correspondence, including
halt. The normalizer's full transition table is materialized and checked;
this is not an unspecified universal-machine index.

\subsection{The finite local-window alphabet}

Use the repository's helical tableau predicate unchanged. A tile records a
vertical-boundary bit, a horizontal-boundary bit, a tape symbol and optional
state, and, on an initialization row, one of the phases L, I, Q, R.
There are exactly
\[
 a_0=2+10(435+1)+4=4366
\]
tiles, with their precise order and local predicate supplied by
\texttt{finite\_window\_constants.py}. An allowed $3\times3$ window is one
whose central cross satisfies that predicate; its four corners are free.
The ordered window list begins with the designated Start and End windows,
then lists all other accepted nine-tuples in lexicographic tile order.

This definition entails at most $4366^9$ tests of a fixed finite predicate.
It is huge but finite and independent of whether K215 halts. A compressed
exact count, with a disjoint-category proof, gives
\[
 k_0=113175624520853754856480928432
\]
windows. The full list is not expanded in this archive. Its definition is
nevertheless complete: there is no unprovided choice of a machine, transition,
tile, window, index, or coefficient.

The local conditions ensure that a Start sits in a finite initialized
rectangle in the periodic covering plane. The phase graph through Q forces
vertical boundaries on both sides, horizontal boundaries propagate across
the strip, and a next horizontal boundary is reached by vertical periodicity.
Between them, the exact update and no-escape rules preserve the unique head.
The final-row head is required to be halted. This is the inherited soundness
argument that excludes spurious accepting cycles. Conversely every finite
one-sided halting run fits into sufficiently large independently padded
width and height, including the powers of five needed by the compiler.

\subsection{The exact fixed integers}

The full coefficient construction is reproduced in
Appendix~\ref{app:coefficient-recipe} and the companion
\texttt{COEFFICIENTS.md}. It specifies every clause coefficient, native
position, copy shift, mask bit, anchor, and optional correction by finite
sums. Its exact interfaces are the reviewed
\texttt{compile\_windows} and modified75 \texttt{new\_constants} recipe,
written out in a standalone form; no historical module is imported at
runtime. In particular:

\begin{enumerate}
\item Each window contributes its selector coefficient and the one-hot
indicators of its nine tile entries. Finitely many coefficient-zero dummy
clauses make room for the native copy positions.
\item Two copies of the center clauses, six horizontal and six temporal
overlap families, and four separated anchors determine sparse polynomials
$D_C,D_R$ in $R_0=2^b$.
\item A single optional high monomial is added to $D_C$ exactly when the
specified finite sum has $2D_C-D_R=0\pmod5$ before that correction. This
bit is a \emph{defined constant}, not a free coefficient or a consistency
test. It ensures the five-adic dummy coefficient is a unit.
\item The native and verification masks are those of the modified75
construction. One ignored extra bit is freed in the native mask and the
low verification bit $4$ is added, retaining the two-bit global population
adjustment.
\item The radix and cell length are the least powers of five meeting the
specified carry and separation bounds. Their evaluated values are
$b=5^{44}$ and $L_0=5^{45}$. Thus $d=5^{89}$.
\end{enumerate}

The integer $d$ itself is
\[
161558713389263217748322010169914619837072677910327911376953125.
\]
Even $B=2^d$ consequently has $d+1$ binary digits. The fact that these
coefficient expansions are infeasible to print must not be confused with
an unknown coefficient. The finite recipe defines them uniquely and proves
their required properties without expanding them. The optional mod-5
correction and all window-dependent sums are retained symbolically.

The delivered controller JSON and coefficient specification are the actual
instantiation artifacts. The archive does not claim to have evaluated every
huge coefficient, enumerated every window, or exhibited a positive solution
to the target equation.

\section{Exact degree and circuit-size verification}

The 70-row evaluator uses only binary additions, subtractions, and
multiplications. A topological and liveness audit checks that every row is
defined before use, every supplied variable and coefficient is live, and
all rows lead to the single output. The count includes all six final
multiplications and the final subtraction of $\Delta$.

For a uniform degree calculation, put
\[
 Q=F+Z+\alpha,\qquad k_0=\eta+\zeta.
\]
These are the leading homogeneous parts of $q$ and $k$.
The exact degrees of $(U_1,U_2,U_0,U_3,U_4,U_5,U_6)$ are
\[
 24,20,1,64,8,2,50,
\]
whose sum is 169. Cancellation in $D^2-\Delta c^2$ must be accounted for;
naively propagating gate degrees gives only an upper bound.
The complete leading homogeneous form is
\begin{equation}\label{eq:leader}
\begin{split}
 -8Q^{104}h\gamma i^4 k_0^{13}w^{13}s^{26}T^2f^2
 \,(Q-\beta J)(wZ-\beta Jt).
\end{split}
\end{equation}
For example, the coefficient of
\[
 J^2F^{104}h\gamma i^4\eta^{13}w^{13}s^{26}T^2f^2t
\]
is $-8\beta^2\ne0$. This proves exact degree 169 for every valid compiler
instance, including the present one. Subtracting $\Delta$, of degree 14,
cannot affect this leading term.

The executable evidence includes exact integer evaluations of the complete
circuits, checks of the scaled-output and packing identities, and independent
dense univariate degree computations at two finite-field specializations.
Those computations supplement the displayed symbolic proof; they are not
the proof of the all-size consistency equivalence or a substitute for the
uniform nonzero leading coefficient.

\section{Nonnegative and unrestricted integer solutions}
\label{sec:domains}

The domain in Theorem~\ref{thm:main} is essential. For a version over
$\N^{16}$, replace each positive unknown $x_j$ by $u_j+1$ and set
\[
 \norms_{\N}(u_1,\ldots,u_{16})
 =\norms_{70}(u_1+1,\ldots,u_{16}+1).
\]
This is a concrete polynomial of the same degree 169, and its absence of
nonnegative solutions is equivalent to $\Con(\mathrm{ZFC})$. The direct
shift adds sixteen additions to the displayed evaluator.

For a polynomial with \emph{unrestricted integer} unknowns, use four
squares for each positive coordinate:
\[
 x_j=1+z_{j,1}^2+z_{j,2}^2+z_{j,3}^2+z_{j,4}^2.
\]
Lagrange's four-square theorem proves that every positive integer occurs.
Hence the explicitly substituted polynomial
\begin{equation}\label{eq:integer-polynomial}
 \norms_{\Z}(\mathbf z)
 =\norms_{70}\left(
 1+\sum_{r=1}^4z_{1,r}^2,\ldots,
 1+\sum_{r=1}^4z_{16,r}^2\right)
\end{equation}
has 64 integer unknowns, exact degree 338, and no integer solution exactly
when ZFC is consistent. To see that the degree is exact, substitute the
sums of four squares into the nonzero degree-169 leading form. Setting
three coordinates in each group to zero leaves its substitution
$x_j=z_{j,1}^2$, which is nonzero because distinct monomials remain
distinct. Hence the degree-338 part does not cancel. The direct evaluator costs $70+16\cdot8=198$
operations. No minimization is claimed for this domain conversion.

Using the raw 16-variable polynomial over all of $\Z^{16}$ would be
incorrect. In particular the proof uses positive $\Delta$ to cancel a
factor, positive ratio slacks, and positive bounds on the packed word.
The four-square substitution supplies those conditions polynomially.
The archive includes the complete 86-row and 198-row evaluators in
\path{artifacts/domain_variants.json}.

\section{Logical and verification status}
\label{sec:logical-status}

The exact Lean endpoint inspected in the pinned repository is
\begin{lstlisting}
BlankHalts machine215 <-> BProv ZFCax_s [] Form.fBot
\end{lstlisting}
and its nonhalting statement is obtained by negating this equivalence.
The repository records completed Lean and Rocq verification for the literal
machine. This work reads those statements and retained records and checks
the table identity; it does not report a new kernel rebuild.

The calibration passes through the compact numerical calculus $C_{222}$:
\[
 \Halt(K215,\mathrm{blank})
 \longleftrightarrow C_{222}(1)
 \longleftrightarrow \mathrm{ZF}\vdash\bot
 \longleftrightarrow \mathrm{ZFC}\vdash\bot.
\]
The last equivalence is the inherited relative-constructibility calibration.
It is about finite refutability, not an assumption that a model of ZFC exists.
The ordinary mathematical proof in the K215 source explicitly identifies
its ZFC formalizability claim. An interpretation of Lean's entire hierarchy
of universes inside ZFC is neither asserted nor needed here.

The present artifact has the following evidence boundaries.

\begin{center}
\begin{tabular}{p{0.34\textwidth}p{0.58\textwidth}}
\toprule
Claim & Evidence supplied\\
\midrule
Exact K215 identity & Pinned table and SHA-256; inherited endpoint statements.\\
One-sided normalization & Explicit table, tape invariant proof, exhaustive
valid folded-row checks, and loader checks.\\
Finite coefficient specification & Exact local predicate, ordered finite
window definition, clause/mask formulas, scalar layout metadata.\\
Fixed-input reduction & Ordinary proof in Sections 4--7, using stated
inherited Pell and compiler lemmas.\\
Arithmetic count and degree & Complete acyclic circuits, independent
formula checks, liveness/count receipts, symbolic leading forms.\\
New formal end-to-end theorem & Not supplied; a proposed formalization task.\\
Expanded giant coefficients or target witness & Not supplied and not
used as evidence.\\
\bottomrule
\end{tabular}
\end{center}

\section{Reproduction and repository integration}

The archive is self-contained for its newly supplied Python checks; only
the standard library is required. From its root run
\begin{lstlisting}
python3 artifacts/fixed_input.py --output artifacts/fixed_input.json
python3 artifacts/domain_variants.py --output artifacts/domain_variants.json
python3 artifacts/compiler/k215_one_sided.py
python3 artifacts/compiler/finite_window_constants.py
python3 artifacts/review/review_fixed70.py
\end{lstlisting}
The root wrapper \texttt{python3 verify.py} runs these commands in order;
the source manifest records the pinned upstream paths and hashes. The routines
perform finite structural and algebraic checks. They do not run K215 until
halt, enumerate its proof search, allocate the full window list, or expand
the giant coefficients.

Compile this article with two runs of pdfLaTeX from the \texttt{report/}
directory. The complete gate appendix is a small generated TeX include,
and is shipped alongside the main source. The root README gives the exact
commands and the generated JSON interfaces.

Suggested integration is a new directory under
\path{Computability/HilbertTenthProblem/Papers/research-wip/}
for the article and artifacts, with cross-links to the pinned K215 endpoint
and the modified75 compiler. The existing universal84 and certificate74
records should remain unchanged: the new statements have fixed-input
contracts. This delivery does not push or modify the user's repository.

\section{Further research questions}

\begin{enumerate}
\item \textbf{Kernel-check the fixed-input adapter.} Formalize the exact
eight-equation projection first, then the auxiliary quotient and signed
repunit reductions. The main new proof obligations are bounded polynomial
identities, positivity, and Lemma~\ref{lem:rep-sign}; the existing K215
endpoint can remain an imported theorem.
\item \textbf{Replace the large finite window enumeration.} The present
recipe is fully explicit but encodes many windows that never occur in a
genuine initialized run. A reachability-preserving sparse atlas could reduce
the coefficient description and height substantially while leaving the
16-variable arithmetic shell unchanged.
\item \textbf{Compile the numerical proof calculus directly.} The repository's
$C_{222}$ calculus has polynomial formula constructors and a short
enumerator. A direct arithmetic compiler might avoid the two-way machine,
folded controller, and local-window expansion. It must still encode
unbounded finite derivations with a fixed number of unknowns; listing
local inference polynomials alone does not solve that problem.
\item \textbf{Optimize with coefficient height charged.} The present
70-operation point treats constants as supplied. A useful second frontier
would measure the bit length or arithmetic description length of every
coefficient and the computational cost of producing it. The best construction
under that metric could be a much larger polynomial with much smaller
coefficients.
\item \textbf{Reduce unknowns without losing the fixed marker.} Sixteen is
not a variable-minimality theorem. A sound elimination of a scale, ratio
slack, or auxiliary coordinate could improve this instance, but any signed
intermediate must be shown positive where the Pell arguments require it.
\item \textbf{Improve the degree/operation tradeoff.} The points
$(70,169)$, $(71,156)$, and the transparent $(84,56)$ describe different
advantages. Can the fixed-input marker and the lower-degree auxiliary
coefficient be combined with a cheaper single-polynomial finalizer?
\item \textbf{Obtain a weaker arithmetic internalization.} The statement
proved here is formalizable in ZFC. A separately formalized internal PA
derivation would require attention to the inherited K215 calibration as
well as to the polynomial compiler. The repository explicitly leaves its
internal PA endpoint distinct from its external proof-assistant theorem.
\end{enumerate}

\appendix
\section{Literal straight-line evaluator}\label{app:circuit}

The following table is generated from the delivered \texttt{fixed70}
packet. The names \texttt{Bm1}, \texttt{Kconstant}, \texttt{MC},
\texttt{MF}, \texttt{W0} denote $\beta,K,M,N,W$, respectively.
In particular \texttt{MF} is already shifted. Every row is one charged
binary operation. Register names are implementation names; the mathematical
definitions are in Section~\ref{sec:polynomial}.

% Generated by build_report.py from artifacts/fixed_input.json.
\begingroup
\footnotesize
\setlength{\tabcolsep}{4pt}
\begin{longtable}{rll}
\caption{The complete 70-operation evaluator.}\\
\toprule
Row & Output register & Binary expression\\
\midrule
\endfirsthead
\toprule
Row & Output register & Binary expression\\
\midrule
\endhead
\midrule
\multicolumn{3}{r}{Continued on the next page}\\
\endfoot
\bottomrule
\endlastfoot
1 & \texttt{\detokenize{C}} & \texttt{\detokenize{Z}} $+$ \texttt{\detokenize{W0}}\\
2 & \texttt{\detokenize{q_minus_F}} & \texttt{\detokenize{C}} $+$ \texttt{\detokenize{alpha}}\\
3 & \texttt{\detokenize{q}} & \texttt{\detokenize{q_minus_F}} $+$ \texttt{\detokenize{F}}\\
4 & \texttt{\detokenize{q_minus_FZ}} & \texttt{\detokenize{q_minus_F}} $-$ \texttt{\detokenize{Z}}\\
5 & \texttt{\detokenize{base_multiple}} & \texttt{\detokenize{Bm1}} $\times$ \texttt{\detokenize{Jrep}}\\
6 & \texttt{\detokenize{norm_repunit}} & \texttt{\detokenize{q}} $-$ \texttt{\detokenize{base_multiple}}\\
7 & \texttt{\detokenize{tau_square}} & \texttt{\detokenize{tau_root}} $\times$ \texttt{\detokenize{tau_root}}\\
8 & \texttt{\detokenize{Lbig}} & \texttt{\detokenize{q}} $\times$ \texttt{\detokenize{q}}\\
9 & \texttt{\detokenize{n2}} & \texttt{\detokenize{Lbig}} $\times$ \texttt{\detokenize{q}}\\
10 & \texttt{\detokenize{wn2}} & \texttt{\detokenize{w}} $\times$ \texttt{\detokenize{Lbig}}\\
11 & \texttt{\detokenize{sn2}} & \texttt{\detokenize{s}} $\times$ \texttt{\detokenize{n2}}\\
12 & \texttt{\detokenize{UM}} & \texttt{\detokenize{wn2}} $\times$ \texttt{\detokenize{sn2}}\\
13 & \texttt{\detokenize{R10b}} & \texttt{\detokenize{eta}} $+$ \texttt{\detokenize{zeta}}\\
14 & \texttt{\detokenize{ksn2}} & \texttt{\detokenize{R10b}} $\times$ \texttt{\detokenize{sn2}}\\
15 & \texttt{\detokenize{first_root_base}} & \texttt{\detokenize{UM}} $\times$ \texttt{\detokenize{ksn2}}\\
16 & \texttt{\detokenize{first_next}} & \texttt{\detokenize{first_root_base}} $+$ \texttt{\detokenize{R10b}}\\
17 & \texttt{\detokenize{first_product}} & \texttt{\detokenize{first_root_base}} $\times$ \texttt{\detokenize{first_next}}\\
18 & \texttt{\detokenize{norm_first}} & \texttt{\detokenize{tau_square}} $-$ \texttt{\detokenize{first_product}}\\
19 & \texttt{\detokenize{R10a}} & \texttt{\detokenize{ksn2}} $+$ \texttt{\detokenize{eta}}\\
20 & \texttt{\detokenize{R12}} & \texttt{\detokenize{UM}} $+$ \texttt{\detokenize{sn2}}\\
21 & \texttt{\detokenize{cam2}} & \texttt{\detokenize{R10a}} $\times$ \texttt{\detokenize{R12}}\\
22 & \texttt{\detokenize{D1}} & \texttt{\detokenize{wn2}} $+$ \texttt{\detokenize{cam2}}\\
23 & \texttt{\detokenize{a4}} & \texttt{\detokenize{4}} $\times$ \texttt{\detokenize{R12}}\\
24 & \texttt{\detokenize{a4m5}} & \texttt{\detokenize{a4}} $+$ \texttt{\detokenize{3}}\\
25 & \texttt{\detokenize{gam}} & \texttt{\detokenize{gamma}} $\times$ \texttt{\detokenize{a4m5}}\\
26 & \texttt{\detokenize{R14}} & \texttt{\detokenize{D1}} $+$ \texttt{\detokenize{gam}}\\
27 & \texttt{\detokenize{L15}} & \texttt{\detokenize{R14}} $\times$ \texttt{\detokenize{R14}}\\
28 & \texttt{\detokenize{a_square}} & \texttt{\detokenize{R12}} $\times$ \texttt{\detokenize{R12}}\\
29 & \texttt{\detokenize{A}} & \texttt{\detokenize{a_square}} $+$ \texttt{\detokenize{a4m5}}\\
30 & \texttt{\detokenize{c2}} & \texttt{\detokenize{R10a}} $\times$ \texttt{\detokenize{R10a}}\\
31 & \texttt{\detokenize{Ac2}} & \texttt{\detokenize{A}} $\times$ \texttt{\detokenize{c2}}\\
32 & \texttt{\detokenize{norm_main}} & \texttt{\detokenize{L15}} $-$ \texttt{\detokenize{Ac2}}\\
33 & \texttt{\detokenize{norm_pair}} & \texttt{\detokenize{norm_first}} $\times$ \texttt{\detokenize{norm_main}}\\
34 & \texttt{\detokenize{norm_triple}} & \texttt{\detokenize{norm_pair}} $\times$ \texttt{\detokenize{norm_repunit}}\\
35 & \texttt{\detokenize{aux_y2}} & \texttt{\detokenize{y_aux}} $\times$ \texttt{\detokenize{y_aux}}\\
36 & \texttt{\detokenize{hpm1}} & \texttt{\detokenize{h}} $\times$ \texttt{\detokenize{UM}}\\
37 & \texttt{\detokenize{index_difference}} & \texttt{\detokenize{R10b}} $-$ \texttt{\detokenize{hpm1}}\\
38 & \texttt{\detokenize{gap_product}} & \texttt{\detokenize{base_multiple}} $\times$ \texttt{\detokenize{q_minus_F}}\\
39 & \texttt{\detokenize{gap}} & \texttt{\detokenize{gap_product}} $+$ \texttt{\detokenize{q_minus_FZ}}\\
40 & \texttt{\detokenize{Lm1}} & \texttt{\detokenize{Lbig}} $-$ \texttt{\detokenize{1}}\\
41 & \texttt{\detokenize{rproduct}} & \texttt{\detokenize{gap}} $\times$ \texttt{\detokenize{Lm1}}\\
42 & \texttt{\detokenize{qMF}} & \texttt{\detokenize{q}} $\times$ \texttt{\detokenize{MF}}\\
43 & \texttt{\detokenize{mask_factor}} & \texttt{\detokenize{MC}} $+$ \texttt{\detokenize{qMF}}\\
44 & \texttt{\detokenize{mask}} & \texttt{\detokenize{mask_factor}} $\times$ \texttt{\detokenize{Jrep}}\\
45 & \texttt{\detokenize{r_lhs}} & \texttt{\detokenize{rproduct}} $+$ \texttt{\detokenize{mask}}\\
46 & \texttt{\detokenize{norm_index}} & \texttt{\detokenize{index_difference}} $-$ \texttt{\detokenize{r_lhs}}\\
47 & \texttt{\detokenize{kinner}} & \texttt{\detokenize{Kconstant}} $+$ \texttt{\detokenize{w}}\\
48 & \texttt{\detokenize{innerC}} & \texttt{\detokenize{kinner}} $\times$ \texttt{\detokenize{C}}\\
49 & \texttt{\detokenize{transport_partial}} & \texttt{\detokenize{innerC}} $+$ \texttt{\detokenize{q_minus_F}}\\
50 & \texttt{\detokenize{local_rhs}} & \texttt{\detokenize{transport_quotient}} $\times$ \texttt{\detokenize{base_multiple}}\\
51 & \texttt{\detokenize{norm_transport}} & \texttt{\detokenize{transport_partial}} $-$ \texttt{\detokenize{local_rhs}}\\
52 & \texttt{\detokenize{L16}} & \texttt{\detokenize{f}} $\times$ \texttt{\detokenize{f}}\\
53 & \texttt{\detokenize{auxiliary_Tf}} & \texttt{\detokenize{auxiliary_quotient}} $\times$ \texttt{\detokenize{f}}\\
54 & \texttt{\detokenize{auxiliary_Tf_minus_one}} & \texttt{\detokenize{auxiliary_Tf}} $-$ \texttt{\detokenize{1}}\\
55 & \texttt{\detokenize{auxiliary_c_Tf}} & \texttt{\detokenize{R10a}} $\times$ \texttt{\detokenize{auxiliary_Tf_minus_one}}\\
56 & \texttt{\detokenize{auxiliary_R_f2}} & \texttt{\detokenize{r_lhs}} $\times$ \texttt{\detokenize{L16}}\\
57 & \texttt{\detokenize{aux_u_rhs}} & \texttt{\detokenize{auxiliary_c_Tf}} $-$ \texttt{\detokenize{auxiliary_R_f2}}\\
58 & \texttt{\detokenize{H2}} & \texttt{\detokenize{aux_u_rhs}} $\times$ \texttt{\detokenize{aux_u_rhs}}\\
59 & \texttt{\detokenize{aux_square_gap}} & \texttt{\detokenize{H2}} $-$ \texttt{\detokenize{aux_y2}}\\
60 & \texttt{\detokenize{aux_coefficient_root}} & \texttt{\detokenize{i}} $\times$ \texttt{\detokenize{Ac2}}\\
61 & \texttt{\detokenize{R16}} & \texttt{\detokenize{aux_coefficient_root}} $\times$ \texttt{\detokenize{aux_coefficient_root}}\\
62 & \texttt{\detokenize{scaled_f_square}} & \texttt{\detokenize{A}} $\times$ \texttt{\detokenize{L16}}\\
63 & \texttt{\detokenize{norm_strong}} & \texttt{\detokenize{scaled_f_square}} $-$ \texttt{\detokenize{R16}}\\
64 & \texttt{\detokenize{L17}} & \texttt{\detokenize{R16}} $\times$ \texttt{\detokenize{aux_square_gap}}\\
65 & \texttt{\detokenize{norm_aux}} & \texttt{\detokenize{L17}} $+$ \texttt{\detokenize{aux_y2}}\\
66 & \texttt{\detokenize{norm_four}} & \texttt{\detokenize{norm_triple}} $\times$ \texttt{\detokenize{norm_aux}}\\
67 & \texttt{\detokenize{norm_product}} & \texttt{\detokenize{norm_four}} $\times$ \texttt{\detokenize{norm_index}}\\
68 & \texttt{\detokenize{all_units}} & \texttt{\detokenize{norm_product}} $\times$ \texttt{\detokenize{norm_transport}}\\
69 & \texttt{\detokenize{seven_units}} & \texttt{\detokenize{all_units}} $\times$ \texttt{\detokenize{norm_strong}}\\
70 & \texttt{\detokenize{polynomial}} & \texttt{\detokenize{seven_units}} $-$ \texttt{\detokenize{A}}\\
\end{longtable}
\endgroup


\section{Complete finite recipe for the five coefficients}
\label{app:coefficient-recipe}

This appendix specifies all five fixed integers in the main polynomial.
Its powers, finite sums, and finite case distinctions concern only fixed
coefficients: none contains a Diophantine unknown. The complete window
list and the binary expansions of the resulting integers are not
materialized. The correction bit defined below is likewise retained as
an exact finite arithmetic expression, rather than claimed to have been
numerically evaluated.

The literal K215 table and its normalized transition function are fixed
in Section~\ref{sec:controller}. For this appendix, denote that normalized
function by
\[
 \delta(q,s)=(q',s',D),\qquad
 q,q'\in\{0,\ldots,434\},\quad
 s,s'\in\{0,\ldots,9\},\quad D\in\{-1,0,1\}.
\]
Its start state is $0$, its erase state is $1$, its absorbing halt is $4$,
and its blank symbol is $0$. In particular $\delta(0,1)=(1,2,0)$.
The archive supplies the entire 4,350-row table, so this use of $\delta$
does not introduce an unspecified universal-machine index.

\subsection{The complete ordered tile alphabet and local predicate}
\label{app:tile-predicate}

Write $\bot$ for an absent payload, phase, or head, as appropriate.
A tile has four fields $(v,h,\pi,\varphi)$: two boundary bits, a payload,
and a phase. A payload is $(s,q)$, with a tape symbol $s$ and a state
$q$ or $\bot$. The ordered tile alphabet is as follows:
\begin{enumerate}
\item Tile $0$ is $\mathsf V=(1,0,\bot,\bot)$ and tile $1$ is
$\mathsf H=(0,1,\bot,\bot)$.
\item For $s=0,\ldots,9$ and $q=\bot,0,\ldots,434$, in this order,
include $(0,0,(s,q),\bot)$. Its tile index is
\[
 2+436s+\eta(q),\qquad
 \eta(\bot)=0,\quad \eta(q)=q+1\quad(q\ne\bot).
\]
\item Tiles $4362,4363,4364,4365$ have phases
$\mathsf L,\mathsf I,\mathsf Q,\mathsf R$, respectively. Their boundary
bits are zero, and their respective payloads are
\[
 (0,\bot),\quad(1,\bot),\quad(1,0),\quad(0,\bot).
\]
\end{enumerate}
Thus the alphabet has $a_c=4366$ elements. All the fields of every tile
are fixed finite data.

Here is the complete cross predicate used to select windows. Let
$t,t_{\ell},t_r,t_d,t_u$ be the center, left, right, down, and up tiles.
Subscripts on fields mean that they are read from the corresponding tile.
First impose
\begin{equation}\label{app:boundary-clause}
 v=v_u,\qquad vv_r=0,\qquad hh_u=0,\qquad
 (v=v_r=0\ \Longrightarrow\ h=h_r).
\end{equation}
If the center is $\mathsf V$ or $\mathsf H$, these are its entire tests.
For a nonboundary center, impose all the following additional conditions.
\begin{enumerate}
\item If $h_d=1$, its phase belongs to
$\{\mathsf L,\mathsf I,\mathsf Q,\mathsf R\}$ and its payload is the
one prescribed above. A left $\mathsf V$ requires phase $\mathsf L$;
a right $\mathsf V$ requires phase $\mathsf R$. If the right neighbor is
not $\mathsf V$, the two phases must form one of the six pairs
\[
 (\mathsf L,\mathsf L),\ (\mathsf L,\mathsf I),\
 (\mathsf I,\mathsf I),\ (\mathsf I,\mathsf Q),\
 (\mathsf Q,\mathsf R),\ (\mathsf R,\mathsf R).
\]
If $h_d=0$, the center phase must instead be $\bot$.
\item If the center carries a head, its transition may not move left
through a left $\mathsf V$ or right through a right $\mathsf V$.
If $h_u=1$, that head, if present, must be in halt state $4$.
\item If $h_u=0$, each non-$\mathsf V$ left or right neighbor must have
a payload. Treat a neighboring $\mathsf V$ as payload $(0,\bot)$.
The next center symbol is its current symbol when headless, and its written
symbol under $\delta$ otherwise. Form the list of incoming states from
the left head moving right, the center head staying, and the right head
moving left. There must be at most one such state. The up payload must
be the next center symbol together with that state, or with $\bot$ if
the list is empty.
\end{enumerate}
This is the literal predicate in the standalone
\path{artifacts/compiler/finite_window_constants.py}; it retains every
initialization, collision, no-escape, and final-row test of the inherited
helical relation~\cite{CompilerDetails}.

A $3\times3$ window is allowed exactly when its central cross satisfies
these tests. Its four corner entries are unrestricted. Write entries in
row-major order, with row $0$ supplying the down neighbor and row $2$
supplying the up neighbor. The two distinguished windows, expressed
entirely in tile indices, are
\[
 \mathsf S=
 \begin{pmatrix}
 1&1&1\\4363&4363&4364\\438&438&876
 \end{pmatrix},\qquad
 \mathsf E=
 \begin{pmatrix}
 1&1&1\\4362&4363&4363\\2&438&438
 \end{pmatrix}.
\]
Both pass the predicate; they are the Start and End windows. The ordered
list $w_0,\ldots,w_{k_c-1}$ begins with $\mathsf S,\mathsf E$ and then
contains every other allowed nine-tuple in lexicographic tile-index order.
This defines a terminating enumeration with at most $4366^9$ predicate
tests, independent of whether the machine halts. Its exact size is
\begin{equation}\label{app:window-count}
 k_c=113175624520853754856480928432.
\end{equation}
The exact compressed census is the sum of
$311471256024587$ allowed central crosses, each with $4366^4$ choices
of corners. It is a finite partition count, not a sample estimate.

\subsection{Evaluated fixed layout parameters}

The following parameters determine the native positions, clause bands,
and radices. The subscript $c$ distinguishes compiler parameters from the
Pell variables in the main polynomial.
\begin{center}
\small
\begin{tabular}{@{}lr@{}}
\toprule
Parameter & Fixed value\\
\midrule
$a_c$ & $4366$\\
$\ell_c$ & $97$\\
$p_c$ & $56587812260426877428240418274$\\
$m_c$ & $113175624520853754856480967730$\\
$U_c$ & $113175624520853754856480993923$\\
$H_c$ & $5771956850563541497680530703172$\\
$T_{c,1}$ & $11883440574689644259930504379381$\\
$T_{c,2}$ & $17994924298815747022180478051224$\\
$g_c$ & $21050666160878798403305464887146$\\
\bottomrule
\end{tabular}
\end{center}
Put
\begin{equation}\label{app:fixed-radices}
 A_c=2^{\ell_c},\qquad b=5^{44},\qquad L_c=5^{45},\qquad
 d=bL_c=5^{89},\qquad R_0=2^b,\qquad B=R_0^{L_c}=2^d.
\end{equation}
These values reproduce the deterministic layout rule:
\begin{align*}
 \ell_c&=\max\{2,\operatorname{bitlength}(k_c+1)\},&
 m_0&=2\ell_c+30a_c+8,\\
 p_c&=\max\left\{0,\left\lceil
       \frac{k_c+9a_c+4-m_0}{2}\right\rceil\right\},&
 m_c&=m_0+2p_c,\\
 U_c&=m_c+6a_c-3,&H_c&=51U_c+3a_c+1,\\
 T_{c,1}&=105U_c+4a_c+2,&T_{c,2}&=159U_c+4a_c+3,\\
 g_c&=186U_c+4a_c+4.
\end{align*}
The least power of five at least $213U_c+4a_c+5$ is $L_c$.
The radix choice is justified after the clause coefficients are defined.

\subsection{Native positions and all clause coefficients}

For $r,s\in\{0,1,2\}$ and $u\in\{0,\ldots,a_c-1\}$, define
\begin{align}
 j_c(r,s,u)&=1+3a_cr+a_cs+u,\label{app:clause-index}\\
 e_c(r,s,u)&=k_c+3a_c+(2-r)3a_c+(2-s)a_c+u.
 \label{app:payload-position}
\end{align}
There is exactly one ignored dummy position,
\[
 e_* = k_c+12a_c =113175624520853754856480980824.
\]
The complete finite native position set is
\begin{align}
 \mathcal E_c={}&\{0,\ldots,k_c-1\}
 \cup\{e_c(r,s,u):0\leq r,s\leq2,\ 0\leq u<a_c\}
 \nonumber\\
 &\cup\{e_*\}\cup\{U_c,3U_c,9U_c,27U_c\}.
 \label{app:native-positions}
\end{align}
Its members are distinct, and $|\mathcal E_c|=m_c+1$.

Define a nonnegative fixed coefficient $c_e$ at each position. For a
selector position $0\leq i<k_c$, set
\begin{equation}\label{app:selector-coefficient}
 c_i=1+
 \sum_{r=0}^{2}\sum_{s=0}^{2}
 A_c^{\,j_c(r,s,(w_i)_{3r+s})}
 +\sum_{h=1}^{4}A_c^{\,9a_c+h}.
\end{equation}
Here $(w_i)_t$ is the tile index at entry $t$ of the already specified
fixed window $w_i$. At the remaining positions put
\begin{align}
 c_{e_c(r,s,u)}&=A_c^{j_c(r,s,u)},&c_{e_*}&=0,\nonumber\\
 c_{U_c}&=A_c^{9a_c+1},&c_{3U_c}&=A_c^{9a_c+2},\nonumber\\
 c_{9U_c}&=A_c^{9a_c+3},&c_{27U_c}&=A_c^{9a_c+4}.
 \label{app:other-coefficients}
\end{align}
The clause mask, including its fixed zero-expression padding clauses, is
\begin{equation}\label{app:clause-mask}
 \mu_c=A_c-2+\sum_{j=1}^{9a_c+4+p_c}A_c^j.
\end{equation}
Omitting the padding summands would change the mask population and would
not be the inherited compiler instance.

For clarity, the least-power-of-five radix rule can now be stated exactly.
Let $C_\Sigma=\sum_{e\in\mathcal E_c}c_e$. Its target is
\[
 T_{\mathrm{radix}}=
 \max\{4(m_c+3)(2C_\Sigma+6)+8,\ 2\mu_c+4,\ 16\}.
\]
Choose $b$ to be the least power of five at least
$\operatorname{bitlength}(T_{\mathrm{radix}}-1)$.
For this instance,
\[
 C_\Sigma<(14k_c+9a_c+5)A_c^{9a_c+4}<2^{3812007},
\]
so the first member of the maximum is below $2^{3812108}$.
The second member dominates. Its exact bit-length contribution is
\begin{align*}
 \operatorname{bitlength}(2\mu_c+3)
 &=\ell_c(9a_c+4+p_c)+2\\
 &=5489017789261407110539324384486.
\end{align*}
Its least power-of-five majorant is $5^{44}$, as asserted
in~\eqref{app:fixed-radices}. No giant power has to be expanded to verify
these inequalities.

\subsection{Convolution coefficients and the fixed correction bit}

The uncorrected center coefficient and the right coefficient are
\begin{align}
 D_{C,0}={}&R_0^{3a_c}+R_0^{H_c+a_c}
             +R_0^{8U_c}+R_0^{24U_c}\nonumber\\
 &+\sum_{e\in\mathcal E_c}c_e
       \bigl(R_0^{T_{c,1}-e}+R_0^{T_{c,2}-e}\bigr),
 \label{app:center-convolution}\\
 D_R={}&R_0^{H_c}.
 \label{app:right-convolution}
\end{align}
All exponents are nonnegative fixed integers. Define the single fixed bit
\begin{equation}\label{app:epsilon}
 \epsilon_c=
 \begin{cases}
 1,&2D_{C,0}-D_R\equiv0\pmod5,\\
 0,&2D_{C,0}-D_R\not\equiv0\pmod5,
 \end{cases}
 \qquad
 D_C=D_{C,0}+\epsilon_c R_0^{g_c}.
\end{equation}
This is a finite arithmetic decision on uniquely defined fixed integers.
It is not a Diophantine variable, an unspecified choice of coefficient,
or a query about consistency. Since $b$ is a power of five,
$R_0\equiv2\pmod5$. The added monomial has nonzero residue after
multiplication by $2$, so the definition guarantees
\[
 2D_C-D_R\not\equiv0\pmod5.
\]
It can in principle be evaluated using modular exponentiation and the
fixed finite sums. This packet does not claim to have evaluated
$\epsilon_c$ or the window-dependent sum in~\eqref{app:center-convolution}.

\subsection{Native masks, shifted field mask, and final ports}

Define the old field mask by the exact finite expression
\begin{align}
 M_{F,0}={}&\mu_c\bigl(R_0^{T_{c,1}}+R_0^{T_{c,2}}\bigr)
 \nonumber\\
 &+\sum_{r=0}^{1}\sum_{s=0}^{2}\sum_{u=0}^{a_c-1}
       R_0^{e_c(r,s,u)}
 \nonumber\\
 &+\sum_{r=0}^{2}\sum_{s=0}^{1}\sum_{u=0}^{a_c-1}
       R_0^{H_c+e_c(r,s,u)}
       +R_0^{9U_c}+R_0^{27U_c}.
 \label{app:old-field-mask}
\end{align}
The modified half-binomial masks are
\begin{align}
 M_C&=B-1-\sum_{\substack{e\in\mathcal E_c\\e\ne1}}R_0^e
           -2R_0^{e_*},\label{app:native-mask}\\
 M_F&=M_{F,0}+4.\label{app:field-mask}
\end{align}
The first modification frees the second bit at the ignored dummy position;
the second adds the low verification bit. These are the modifications of
the reviewed compiler export~\cite{CompilerDetails}. In particular
\begin{gather*}
 0<M_C,M_F<B-1,\qquad
 M_C\equiv2\pmod4,\qquad M_F\equiv4\pmod8,\\
 \pop(M_C)+\pop(M_F)=d.
\end{gather*}

The five coefficients supplied to the main polynomial are now exactly
\begin{equation}\label{app:five-final-coefficients}
 \boxed{
 \begin{aligned}
 \beta&=B-1,&K&=D_C+BD_R,\\
 M&=M_C,&N&=M_F+B-1,\\
 W&=R_0B^2=2^{\,2\cdot5^{89}+5^{44}}.
 \end{aligned}}
\end{equation}
The source coefficient $N$ is already shifted. It is \emph{not} the
native one-cell mask $M_F$; replacing one by the other changes the
polynomial. The endpoint $W$ inserts the End selector at cyclic cell $2$,
giving the three-I initialization run and four unary ones already proved
equivalent to blank K215.

\subsection{Explicitness and the representation boundary}

Every summation bound, finite set, field lookup, and exponent in this
appendix is fixed before the sixteen unknowns of the main polynomial
are chosen. Substitution of~\eqref{app:five-final-coefficients} therefore
produces one ordinary integer-coefficient polynomial. No variable power,
bounded universal quantifier, sequence parameter, or acceptance test
remains in that polynomial.

The compressed representation is essential for practical delivery:
$B$ alone has $5^{89}+1$ binary digits. This coefficient height is not
included in the 70-operation evaluation count. The count assumes access
to the five fixed numerals, as stated in the main theorem; it does not
assert that their construction or storage is inexpensive. The companion
\path{artifacts/compiler/COEFFICIENTS.md}, source code, and pinned
manifests give the same recipe in machine-readable and prose forms.

All compiler formulas in this appendix are the lazy form of the retained
construction
\[
 \texttt{new\_constants(compile\_windows(windows,a))}.
\]
The new explicit controller instantiation does not replace the inherited
proof of that compiler, and the exact finite specification must not be
confused with a claim that all its giant integers have been expanded.


\section{Source and dependency references}

The source pin for every ProveIt item below is \repo. The delivered
provenance manifest records full paths and byte hashes for the source
artifacts used. These references distinguish inherited theorems from the
new fixed-input reductions.

\begin{thebibliography}{20}
\bibitem{Jones1980}
J. P. Jones, \emph{Undecidable Diophantine equations}, Bulletin of the
American Mathematical Society (N.S.) 3 (1980), 859--862.
\url{https://doi.org/10.1090/S0273-0979-1980-14832-6}.

\bibitem{Jones1982}
J. P. Jones, \emph{Universal Diophantine equation}, Journal of Symbolic
Logic 47 (1982), 549--571.
\url{https://www.cambridge.org/core/journals/journal-of-symbolic-logic/article/113EBF3D928DFC6011C52BECAE81A087}.

\bibitem{CarlMoroz}
M. Carl and B. Z. Moroz, \emph{On a Diophantine representation of the
predicate of provability}, Zap. Nauchn. Sem. POMI 407 (2012), 77--104;
English translation, Journal of Mathematical Sciences 199 (2014), 36--52.
\url{https://www.mathnet.ru/eng/znsl5486}.

\bibitem{FormalDPRM}
D. Larchey-Wendling and Y. Forster, \emph{Hilbert's Tenth Problem in Coq},
Logical Methods in Computer Science 18(1) (2022).
\url{https://arxiv.org/abs/2003.04604}.

\bibitem{K215}
ProveIt, \emph{Formal correctness of the exact 215-state machine};
\path{Computability/BusyBeaver/Research/ZFC/k225-formal/verification/k215/README.md};
ordinary proof in \path{candidate-k215/proof.md}; endpoint in
\path{Computability/BusyBeaver/Lean/BusyBeaver/ZFC/KEndpoint215.lean}.
\url{https://github.com/VladimirReshetnikov/ProveIt/tree/cdd582bec13f908df3711fa9e05d1ddac368a5db/Computability/BusyBeaver/Research/ZFC/k225-formal/verification/k215}.

\bibitem{Kernel}
ProveIt, \emph{A strong half-binomial Pell kernel in 42 operations},
\path{Computability/HilbertTenthProblem/Papers/research-wip/native-stream-queue/pell_kernel_half_binomial42.md}.
Its references include the relaxed auxiliary-rank and half-parameter
signed step-down proofs.

\bibitem{Asymmetric}
ProveIt, \emph{Asymmetric scales lower the complete87/88 degrees},
\path{Computability/HilbertTenthProblem/Papers/research-wip/native-stream-queue/complete75_asymmetric_scale_tradeoffs.md},
especially Sections 2--4.

\bibitem{Bezout}
ProveIt, \emph{An auxiliary quotient gives a complete universal polynomial
in 85 operations},
\path{Computability/HilbertTenthProblem/Papers/research-wip/native-stream-queue/complete85_auxiliary_bezout_projection.md},
especially Sections 3--5.

\bibitem{Compiler75}
ProveIt, \emph{A 75-operation fixed-index universal certificate},
\path{Computability/HilbertTenthProblem/Papers/1980/FIXED_RAW_UNIVERSAL_75_PROOF.md}.

\bibitem{CompilerDetails}
ProveIt, \emph{The modified half-binomial compiler},
\path{Computability/HilbertTenthProblem/Papers/research-wip/native-stream-queue/complete75_half_binomial_compiler.md},
with its executable layout recipe and the preceding fixed76 compiler,
helical tableau, and cyclic-marker proofs.

\bibitem{Universal84}
ProveIt, \emph{A complete universal polynomial in 84 operations},
\path{Computability/HilbertTenthProblem/Papers/research-wip/native-stream-queue/complete84_scaled_strong_output.md},
with the literal 84-row JSON source and its independent review.

\bibitem{Certificate74}
ProveIt, \emph{A factored first norm lowers the complete comparison bound
to 74},
\path{Computability/HilbertTenthProblem/Papers/research-wip/native-stream-queue/complete74_factored_first_norm.md}.
\end{thebibliography}

\end{document}
